\cotefile{157-2}
\scanpage{307}

\mspage{48}
% archive sheet number: 306

6. \uline{\unsure{Compositions|Composition} des} $\underline{\mathrm{Ty}}$

On a une loi de composition (partout définie) dans $\operatorname{Ob}(\underline{\mathrm{Ty}})=\mathrm{Ty}$, c'est celle du monoïde libre à deux générateurs $\Delta_+,\Delta_-$. Cette loi de composition s'interprète, au niveau \deleted{des intervalles} catégorique (et non plus \unsure{seulement} combinatoire) comme une somme amalgamée
\[
\begin{aligned}
&\deleted{I_\tau I_{\tau'}=\unsure{I_{\tau,\tau'}|I_{\tau\tau'}}}\\
&I_{\tau\tau'}=I_\tau\amalg_{(e,\underline{1}_\tau,\underline{0}_{\tau'})}I_{\tau'}
\end{aligned}
\]
où $e$ est la catégorie unité, et où $\underline{1}_\tau,\underline{0}_{\tau'}$ représentent \unsure{maintenant} les morphismes $e\unsure{\text{ --- }|\longrightarrow}I_\tau$, $e\longrightarrow I_{\tau'}$ ayant comme image l'élément but et l'élément source respectivement de l'intervalle orienté envisagé. On a \unsure{donc} un carré cocartésien dans $(\mathrm{Cat})$
\diagram{\begin{tikzcd}[ampersand replacement=\&,column sep=small,row sep=small]
\& e \arrow[dl,hook,"\underline{1}_\tau"'] \arrow[dr,hook,"\underline{0}_{\tau'}"] \& \\
I_\tau \arrow[dr,hook,"i_{\tau,\tau'}"'] \& \text{\unsure{cocart}} \& I_{\tau'} \arrow[dl,hook,"j_{\tau,\tau'}"] \\
\& I_{\tau\tau'} \& {}
\end{tikzcd}},
où toutes les flèches sont des foncteurs mono et pleinement fidèles. Que \unsure{le} diagramme \inserted{(évidemment commutatif)} soit bien cocartésien \inserted{dans $(\mathrm{Cat})$} revient, par définition, à dire que pour \deleted{\unsure{un}} toute $X$ dans $\mathrm{Cat}$, le diagramme transposé
\diagram{\begin{tikzcd}[ampersand replacement=\&,column sep=small,row sep=small]
\& \deleted{\unsure{\mathrm{Ch}_e|\mathrm{Ch}}}(X) \& \\
\mathrm{Ch}_\tau(X) \arrow[ur,phantom,"\deleted{\swarrow}"{description}] \arrow[ur,"\beta_\tau"{description},"\text{(but)}"'] \& \& \mathrm{Ch}_{\tau'}(X) \arrow[ul,phantom,"\deleted{\searrow}"{description}] \arrow[ul,"\leftarrow\!\sigma_{\tau'}\;\text{(source)}"']\\
\& \mathrm{Ch}_{\tau\tau'}(\unsure{X|\mathcal X}) \arrow[ul,phantom,"\deleted{\searrow}"{description}] \arrow[ur,phantom,"\deleted{\swarrow}"{description}] \arrow[ul] \arrow[ur] \&
\end{tikzcd}}
\scanpage{308}

\mspage{49}
% archive sheet number: 307

ce qui \deleted{signifie} revient à la tautologie suivante qu'un chemin \inserted{$C$} dans $X$ de type $\tau\tau'$ ``n'est autre chose'' qu'un couple $(c,c')$ d'un chemin \inserted{$c$} de type $\tau$ et d'un chemin $c'$ de type $\tau'$ ``mis bout à bout'' (dans cet ordre), i.e. tels que
\[
\operatorname{but}(c)=\operatorname{source}(c')
\]
(et on a alors
\[
\begin{aligned}
\operatorname{\unsure{source}}(C)&=\operatorname{source}(c)\\
\operatorname{but}(C)&=\operatorname{but}(c')\qquad ).
\end{aligned}
\]

Il se pose ici le sempiternel pb si on va noter ce chemin $c.c'$, ce qui semble à première vue s'imposer, ou plutôt $c'\circ c$ (\deleted{\unsure{en}} \unsure{s'imaginant|imaginant} \deleted{$c$}\inserted{le chemin $c$ de $x=\operatorname{source}(c)$ vers $y=\operatorname{but}(c)$} comme une \deleted{application} sorte de morphisme de $x$ \deleted{\unsure{vers}} $y$ \deleted{\unsure{et}} (cf plus bas le formalisme des cat de fractions), $c'$ comme un morphisme de $y=\deleted{\unsure{so}}\operatorname{source}(c')$ vers $z=\operatorname{but}(c')$, et en \unsure{interprétant} le composé $C$ comme le composé de ces deux morphismes \inserted{$c,c'$} (dans cet ordre du point de vue des opérations à effectuer), qu'on note traditionnellement $c'\circ c$ ou $c'.c$, et non $c\circ c'$ ou $c.c'$, hélas! C'est cette notation que nous allons suivre plus bas.

Mais pour l'instant on ne va pas se soucier des chemins dans des cat. quelconques $X$, mais revenir au fonctorialisme des $\tau$ et des $I_\tau$ eux-mêmes. L'interprétation catégorique de $\unsure{\tau\tau'|I_{\tau'}}$ comme une somme amalgamée nous assure que cette opération de composition est fonctorielle :
\scanpage{309}

% archive sheet number: 308

\[
\tau_1\xrightarrow{f_{\deleted{\unsure{1}}}}\tau_2,\quad
\tau'_1\xrightarrow{f'_{\deleted{\unsure{2}}}}\tau'_2
\qquad\text{donne}\qquad
\tau_1\tau'_1\xrightarrow{f_{\deleted{\unsure{1}}}*f'_{\deleted{\unsure{2}}}}\tau_2\tau'_2
\]
(on \unsure{n'utilisera|utilisera} pas $f_{\deleted{\unsure{1}}}\overset{\unsure{*}}{\unsure{\cdot|\circ}}f'_{\deleted{\unsure{2}}}$ pour la flèche notée $f_{\deleted{\unsure{1}}}*f'_{\deleted{\unsure{2}}}$, à cause de la confusion possible avec la composition des flèches dans $\underline{\mathrm{Ty}}$) : On a transitivité de façon évidente si de plus $\tau_2\overset{g_{\deleted{\unsure{1}}}}{\unsure{\text{---}|\longrightarrow}}\tau_3$, $\tau'_2\overset{g'}{\unsure{\text{---}|\longrightarrow}}\tau'_3$
\[
\begin{aligned}
&\deleted{\unsure{(g_1f_1)*(g_2|(gf)*(g'|(gf)*g'}}\\
&(gf)*(g'f')=(g*g')\circ(f*f').
\end{aligned}
\]

Donc on a un bifoncteur de composition
\[
\underline{\mathrm{Ty}}\times\underline{\mathrm{Ty}}\longrightarrow\underline{\mathrm{Ty}}
\]
\margin{$(c,c')\longmapsto c.c'$\\ ou $c*c'$}

(En fait, \deleted{\illegible{2}}\inserted{on a aussi} un bifoncteur \deleted{$\unsure{\mathrm{Ty}|\mathrm{Typ}}$}
\[
\underline{\mathrm{Typ}}\times\underline{\mathrm{Typ}}\unsure{\text{ --- }|\longrightarrow}\underline{\mathrm{Typ}},
\]
mais on préfère oublier maintenant la catégorie $\underline{\mathrm{Typ}}$, qui donne un formalisme nettement moins simple.)

Cette composition est \uline{strictement} associative et \uline{strictement} \uline{unitaire} (l'unité étant bien sûr l'élément unité $\unsure{1|s}$ du monoïde libre $\mathrm{Ty}$, correspondant à la catégorie ponctuelle $\{0\}$), et pas seulement ``à isom près''. Et pour cause, puisque dans $\underline{\mathrm{Ty}}$ les \deleted{\unsure{objets}} iso sont des identités. \uline{Donc} $\underline{\mathrm{Ty}}$ \uline{est un monoïde dans la catégorie} $(\mathrm{Cat})$ des petites catégories!
\scanpage{310}

\mspage{49}
% archive sheet number: 309

Cela implique que non seulement $\operatorname{Ob}\underline{\mathrm{Ty}}=\mathrm{Ty}$, mais aussi $\operatorname{Fl}(\underline{\mathrm{Ty}})$, et plus généralement tous les ensembles \deleted{\unsure{de|tels}} de la forme $\operatorname{Hom}_{\mathrm{Cat}}(X,\underline{\mathrm{Ty}})$, \inserted{sont des monoïdes, et ceci fonctoriellement en $X$ ;} par exemple les composantes $\operatorname{Hom}_{\mathrm{Cat}}(\Delta_i,\underline{\mathrm{Ty}})$ de l'ens. semi-simplicial ``nerf'' de $\underline{\mathrm{Ty}}$
\[
\operatorname{Nerf}(\underline{\mathrm{Ty}})_i=\operatorname{Hom}_{\mathrm{Cat}}(\Delta_{\unsure{i}},\underline{\mathrm{Ty}}),
\]
sont des monoïdes, et les applications \unsure{sesimpl.|ssimpl.} entre elles sont des hom. de monoïdes : le \unsure{Nerf} \deleted{est} de $\underline{\mathrm{Ty}}$ est un \uline{monoïde} \uline{semi-simplicial}.

On va décrire ces monoïdes $\operatorname{Nerf}_i(\underline{\mathrm{Ty}})$, et constater que ce sont \inserted{tous} des monoïdes \uline{libres}, \deleted{\unsure{et}} \deleted{\unsure{explicit.}} tout comme $\operatorname{Nerf}_0(\underline{\mathrm{Ty}})=\operatorname{Ob}\underline{\mathrm{Ty}}=\mathrm{Ty}$. Les applications \unsure{sspl.} entre ces monoïdes sont donc connues quand on les connaît sur les \uline{générateurs} de ces monoïdes libres, (\deleted{\unsure{c. à d. sur les}}\inserted{lesquels générateurs} sont aussi les éléments indécomposables, donc définis intrinsèquement en termes de \unsure{sa|la} structure multiplicative).

Pour $\operatorname{Nerf}_0(\underline{\mathrm{Ty}})=\mathrm{Ty}$, c'est une chose vue, les générateurs sont
\[
\Delta_+,\Delta_-,
\]
correspondant respectivement aux intervalles orientés
\diagram{\begin{tikzcd}[ampersand replacement=\&,row sep=0pt,column sep=small]
\Delta_+: \& 0 \arrow[r] \& 1\\
\Delta_-: \& 0 \& 1 \arrow[l]
\end{tikzcd}.}
\scanpage{311}

\mspage{50}\mspage{50}
% archive sheet number: 310

Pour \inserted{$\operatorname{Nerf}_1(\underline{\mathrm{Ty}})=\operatorname{Fl}(\underline{\mathrm{Ty}})$, i.e. pour} les flèches, les résultats \deleted{\unsure{du}} \unsure{]} du par. précédent (détermination de $\operatorname{Fl}(\underline{\mathrm{Ty}})$) s'interprètent aussi en disant que $\operatorname{Fl}(\underline{\mathrm{Ty}})$ est le monoïde libre engendré par un système de générateurs infini
\[
f_{\varepsilon,n},\qquad(\varepsilon,n)\in\{+1,-1\}\times\mathbb N
\]
indexé par \unsure{l'ens produit|leur produit|un produit} \deleted{\unsure{que}} $\{+1,-1\}\times\unsure{\mathcal N|\mathbb N}$. Ici $f_{\varepsilon,n}$ désigne la \inserted{unique} flèche de $\mathrm{Ty}$
\[
f_{\varepsilon,\ell}:\Delta_\varepsilon\longrightarrow\Delta_\varepsilon^{\deleted{\unsure{n}}\ell}
\]
(où bien sûr l'exposant $n$ désigne une exponentiation pour la loi de composition de $\underline{\mathrm{Ty}}$
\[
\Delta_\varepsilon^\ell=\underbrace{\Delta_\varepsilon\circ\Delta_\varepsilon\circ\cdots\circ\Delta_\varepsilon}_{\ell\ \text{copies}},
\]
et non la puissance cartésienne dans $\mathrm{Cat}$, ou dans $\underline{\mathrm{Ty}}$ (où elle n'existe que si $n\in\{0,1\}$)).

\deleted{\unsure{on}} Dans la description de $\operatorname{Fl}(\underline{\mathrm{Ty}})$ dans le par. précédent, on a écrit le ``terme général'' du monoïde libre en groupant ensemble les \deleted{\unsure{f}}\inserted{groupes de} facteurs successifs de la forme $f_{+,0}, f_{-,0}$ i.e. correspondant \deleted{\unsure{aux deux}}\inserted{aux} flèches \inserted{de} dégénérescence indécomposables (pour $*$)
\[
\left\{\begin{aligned}
 f_{+,0}:{}&\Delta_+\unsure{\text{ --- }|\longrightarrow}e\\
 f_{-,0}:{}&\Delta_-\unsure{\text{ --- }|\longrightarrow}e
\end{aligned}\right.
\]
% slanted
\margin{\uline{NB} Le ``mot'' vide de $\operatorname{Fl}(\underline{\mathrm{Ty}})$ \deleted{\unsure{est}} correspond à la flèche $\mathrm{id}_1$, où $1$ est l'élément unité de $\mathrm{Ty}$, \deleted{\unsure{cor}} \unsure{représenté} par l'``intervalle'' ponctuel}
Le \deleted{\unsure{grou}} sous-\deleted{groupe}monoïde libre de $\operatorname{Fl}(\underline{\mathrm{Ty}})$ engendré par ces deux éléments \deleted{s'identifie \unsure{avec}}\inserted{est le} monoïde des flèches de dégénérescence \inserted{(pures)} quelconques,
\scanpage{312}

\mspage{51}
% archive sheet number: 311

de $\underline{\mathrm{Ty}}$ :
\[
f_\tau:I_\tau\unsure{\text{ --- }|\longrightarrow}e,\qquad(\tau\in\mathrm{Ty})
\]
et on a
\[
f_\tau*f_{\tau'}=f_{\tau\tau'},
\]
i.e. \deleted{\unsure{la multip.}} l'application
\diagram{\begin{tikzcd}[ampersand replacement=\&,row sep=small,column sep=small]
\tau \arrow[r,mapsto] \& f_\tau : \& \operatorname{Ob}(\underline{\mathrm{Ty}}) \arrow[r] \arrow[d,equal] \& \operatorname{Fl}(\underline{\mathrm{Ty}}) \arrow[dd,equal]\\
\& \& \mathrm{Ty} \arrow[d,equal] \& \deleted{\unsure{\mathrm{Ty}}}\\
\& \& \operatorname{Nerf}_0(\underline{\mathrm{Ty}}) \& \operatorname{Nerf}_1(\underline{\mathrm{Ty}})
\end{tikzcd}}
est un homomorphisme de monoïdes, transformant $\Delta_\varepsilon$ $(\varepsilon\in\{+1,-1\})$ dans $\deleted{\unsure{\varepsilon}} f_{\varepsilon,0}$. \deleted{Elle n'a pas de \unsure{mérite|intérêt}, puisque c'est l'application de} \deleted{L'application \unsure{de|se}} Ce n'est pas l'application de dégénérescence de $\operatorname{Nerf}_0$ dans $\operatorname{Nerf}_1$, i.e. l'application $\tau\unsure{\rightarrow|\text{ --- }}\mathrm{id}_\tau$. Celle ci est donnée par
\[
\begin{aligned}
 \operatorname{Ob}(\underline{\mathrm{Ty}})&\longrightarrow \operatorname{Fl}(\mathrm{Ty})\\
 \tau&\mathrel{\mapstochar\relbar}\mathrm{id}_\tau
\end{aligned}
:\quad
\left|\begin{aligned}
\Delta_\varepsilon&\longmapsto f_{\varepsilon,1}\qquad(\varepsilon\in\{+1,-1\})\\
\text{i.e.}\qquad \Delta_+&\mathrel{\mapstochar\relbar} f_{+,1}\\
\Delta_-&\mathrel{\mapstochar\relbar} f_{-,1}.
\end{aligned}\right.
\]

Quant aux homomorphismes source et but
\diagram{\begin{tikzcd}[ampersand replacement=\&,column sep=large]
\operatorname{Fl}(\underline{\mathrm{Ty}})\arrow[r,shift left,"\sigma\ (\text{source})"]\arrow[r,shift right,"\beta\ (\text{but})"']\&\operatorname{Ob}(\underline{\mathrm{Ty}})
\end{tikzcd},}
elles sont \unsure{évidentes} sur les générateurs
\[
\left\{\begin{aligned}
\sigma(f_{\varepsilon,\ell})&=\Delta_\varepsilon\\
\beta(f_{\varepsilon,\ell})&=\Delta_\varepsilon^\ell.
\end{aligned}\right.
\]

Je vais décrire maintenant $\operatorname{Nerf}_2(\mathrm{Ty})$, i.e. les \deleted{\unsure{systè}} systèmes
\diagram{\begin{tikzcd}[ampersand replacement=\&,column sep=normal]
\tau\arrow[r,"f"] \& \tau'\arrow[r,"g"] \& \tau''.
\end{tikzcd}}
\scanpage{313}

\mspage{52}
% archive sheet number: 312

Se donner $f\in\operatorname{Fl}(\underline{\mathrm{Ty}})$ revient à se donner une suite ordonnée de générateurs $f_{\varepsilon\,\unsure{;|,}\ell}$, donc une suite
\[
\begin{aligned}
&\deleted{\unsure{f_{(\varepsilon_1,\ell_1)}|f_{(\varepsilon_i,\ell_i)}}}\\
&f_{\varepsilon_1,\ell_1},\quad f_{\varepsilon_2,\ell_2}\quad\cdots\quad f_{\varepsilon_n,\ell_n}\\
&\big((\varepsilon_i,\ell_i)\big)_{1\leq i\leq n}\in[\{\pm1\}\times\mathbb N]^n.
\end{aligned}
\]

On a alors
\[
\begin{aligned}
x&=\Delta_{\varepsilon_1}^{\deleted{\unsure{\ell_1}}}\Delta_{\varepsilon_2}^{\deleted{\unsure{\ell_2}}}\cdots\Delta_{\varepsilon_n}^{\deleted{\unsure{\ell_n}}}&\operatorname{long}x&=n\unsure{.}\\
y&=\Delta_{\varepsilon_1}^{\ell_1}\Delta_{\varepsilon_2}^{\ell_2}\cdots\Delta_{\varepsilon_n}^{\ell_n}&\operatorname{long}y&=\sum_{i\in[1,n]}\ell_i
\end{aligned}
\]

Se donner $g:y\unsure{\text{ --- }|\longrightarrow}z$, i.e. un objet de $y\backslash\underline{\mathrm{Ty}}\deleted{\unsure{/z}}$, revient à se donner une application de l'ens des $\deleted{\unsure{n=}}m=\sum\ell_i$ flèches de transition de $y$ dans $\mathbb N$. En fait, on a
\[
\operatorname{Tr}(y)=\operatorname{Tr}(\Delta_{\varepsilon_1}^{\ell_1})\amalg
\operatorname{Tr}(\Delta_{\varepsilon_2}^{\ell_2})\cdots\amalg\operatorname{Tr}(\Delta_{\varepsilon_n}^{\ell_n})
\]
(\inserted{``une somme'', i.e.} réunion disjointe), donc
\[
\operatorname{Hom}_{\mathrm{Ens}}\unsure{(\operatorname{Tr}y))|(\operatorname{Tr}y)}
\simeq\prod_{i\in[1,\unsure{n|m}]}\operatorname{Hom}_{\mathrm{Ens}}
 \big(\underbrace{\operatorname{Tr}(\Delta_{\varepsilon_i}^{\ell_i})}_{[1,\ell_i]},\mathbb N\big)
\]
donc la donnée de $g$ équivaut à la donnée d'un système de $n$ applications
\[
\left\{\begin{aligned}
\pi_i:[1,\ell_i]&\longrightarrow\mathbb N\qquad(\text{pondération de }\Delta_{\varepsilon_i}^{\ell_i})\\
&\text{pour}\quad i\in[1,n]
\end{aligned}\right.
\]

Mais pour $i\in[1,n]$, considérons \unsure{l'|la} \deleted{\unsure{chaîne}}\inserted{\unsure{élément de|chaîne de|chemin de}} \deleted{flèche} $\operatorname{Nerf}_2$
\diagram{\[
\begin{tikzcd}[ampersand replacement=\&,column sep=normal]
\Delta_{\varepsilon_i}\arrow[r,"{f_{\varepsilon_i,\unsure{\ell_i|c_i}}}"] \&
\Delta_{\varepsilon_i}^{\ell_i}\arrow[r,"{g_{\unsure{\varepsilon_i|\deleted{g}},\pi_i}}"] \&
\underbrace{\Delta_{\varepsilon_i}^{\pi_i(1)}\Delta_{\varepsilon_i}^{\pi_i(\unsure{1|2})}\cdots
\Delta_{\varepsilon_i}^{\pi_i(\ell_i)}}_{\Delta_{\varepsilon_i}^{\sum_j\pi_i(j)\,\overset{\mathrm{déf}}{=}\operatorname{long}(\pi_i)}}
\end{tikzcd}
\tag*{(1)}
\]}
\scanpage{314}

\mspage{53}
% archive sheet number: 313

où $g_{\varepsilon_i,\pi_i}$ est la flèche définie en \unsure{terme|termes} de $\pi_i$ au par. précédent. Mais \unsure{ce n'est autre aussi|on n'en note aussi} que
\[
g_{\varepsilon_i,\deleted{\unsure{\ell_i}}\pi_i}
=f_{\varepsilon_i,\pi_i(1)}*f_{\varepsilon_i,\pi_i(\unsure{2|1})}*
\cdots*f_{\varepsilon_i,\pi_i(\ell_i)}.
\tag*{(2)}
\]

L'élément de $\operatorname{Nerf}_2$ envisagé ne dépend que de
\[
(\varepsilon_i,\ell_i,\pi_i),\ \text{où}\quad
(\varepsilon_i,\ell_i)\in\{\pm1\}\times\mathbb N,\quad
\pi_i:[1,\ell_i]\longrightarrow\mathbb N.
\]
(\uline{NB} si $\ell_i=0$, $[1,\ell_i]=[1,0]=\varnothing$ par convention, et $\pi_i$ est l'unique application $\varnothing\to\mathbb N$, \unsure{aussi|donc} la donnée $(\varepsilon_i,\ell_i,\pi_i)$ revient alors à $(\varepsilon_i,\ell_i)$\unsure{.})

Si
\[
(\varepsilon,\ell,\pi)\quad\text{comme dessus},\qquad
(\varepsilon,\ell)\in\{\pm1\}\times\mathbb N,\quad\pi:[\unsure{1|0},\ell]\to\mathbb N,
\]
désignons par
\[
f_{\varepsilon,\ell,\pi}
\]
l'élément de $\operatorname{Nerf}_2$ donné par
\diagram{\begin{tikzcd}[ampersand replacement=\&,column sep=normal]
\Delta_\varepsilon\arrow[r,"{f_{\varepsilon,\ell}}"] \&
\Delta_\varepsilon^\ell\arrow[r,"{f_{\varepsilon,\pi(1)}*\cdots*f_{\varepsilon,\pi(\ell)}}"] \&[5em]
\Delta_\varepsilon^{\pi(1)}\Delta_\varepsilon^{\pi(\unsure{1|2})}\cdots\Delta_\varepsilon^{\pi(\ell)}
\end{tikzcd}}

On peut \deleted{\unsure{le}} formuler la construction précédente en disant que tout élément de $\operatorname{Nerf}_2$ peut se décomposer de façon unique en un produit \unsure{*}
\[
(f,g)=f_{\varepsilon_1,\ell_1,\pi_1}*f_{\varepsilon_2,\ell_2,\pi_2}*
\cdots*f_{\varepsilon_n,\ell_n,\pi_n},
\]
ce qui prouve que $\operatorname{Nerf}_2(\underline{\mathrm{Ty}})$ est le monoïde libre sur les générateurs précédents $f_{\varepsilon,\ell,\pi}$.

(\uline{NB} : la suite vide correspond
\diagram{\begin{tikzcd}[ampersand replacement=\&,column sep=normal]
1\arrow[r,"\mathrm{id}_1"] \&
1\arrow[r,phantom,"{\overset{\mathrm{id}_1}{\unsure{\text{ --- }|\longrightarrow}}}"{description}] \& 1\quad .)
\end{tikzcd}}
\scanpage{315}

\mspage{54}
% archive sheet number: 314

Si l'objet de $\operatorname{Nerf}_2$ est
\diagram{\begin{tikzcd}[ampersand replacement=\&,column sep=normal]
x\arrow[r,phantom,"{\overset{f}{\unsure{\text{ --- }|\longrightarrow}}}"{description}] \&
y\arrow[r,dash,"g"] \& z,
\end{tikzcd}}
la construction précédente nous dit, en termes de la donnée combinatoire i.e. de la suite de \deleted{\unsure{données}}
\[
(\varepsilon_1,\ell_1,\pi_1),\ldots,(\varepsilon_n,\ell_n,\pi_n),
\]
quels sont $x,y,z$ : chacun d'eux est le produit dans $\deleted{\unsure{\operatorname{Ob}}}\operatorname{Nerf}_0=\mathrm{Ty}$ des $n$ facteurs $x_i,y_i$\inserted{\unsure{ou|et|,}} $z_i$ respectivement, pour $i\in[1,n]$, donnés dans la définition des $f_{(\varepsilon_i,\ell_i,\pi_i)}$, i.e. de
\[
\Delta_{\varepsilon_i},\quad\Delta_{\varepsilon_i}^{\ell_i},\quad\text{et}\quad
\Delta_{\varepsilon_i}^{\pi_i(1)}\Delta_{\varepsilon_i}^{\pi_i(\unsure{1|2})}\cdots
\Delta_{\varepsilon_i}^{\pi_i(\ell_i)}=\Delta_{\varepsilon_i}^{\operatorname{long}(\pi_i)}
\]
(où $\operatorname{long}\pi_i=\sum_{j\in[1,\ell_i}\pi_i(j)$). Donc
\[
\left\{\begin{aligned}
x&=\prod_{i\in[1,n]}^{\mathrm{Ty}}\Delta_{\varepsilon_i}\\
y&=\prod_{i\in[1,n]}^{\mathrm{Ty}}\Delta_{\varepsilon_i}^{\ell_i}
\deleted{\unsure{=\Delta_{\varepsilon_i}^{\operatorname{long}(\ell)}}}\\
\unsure{Z|z}&=\prod_{i\in[1,n]}^{\mathrm{Ty}}\Delta_{\varepsilon_i}^{\operatorname{long}(\pi_i)}
\deleted{\unsure{=\Delta_{\varepsilon_i}^{\sum_{i\in[1,n]}\operatorname{long}(\pi_i)}}},
\end{aligned}\right.
\]
\margin{attention, les produits doivent être pris en respectant l'ordre des facteurs donné par la \unsure{notation} indicielle en $i$}
et on trouve de même $f,g$ comme des produits (au sens $*$) de $n$ facteurs $f_i$, resp. $g_i$ $(i\in[1,n])$ explicités dans la formule (1)(2) (p. 52, 53) :
\[
\begin{aligned}
\unsure{\{|[} \qquad f_i&=f_{\varepsilon_i,\ell_i}:\Delta_{\varepsilon_i}\longrightarrow\Delta_{\varepsilon_i}^{\ell_i}\\
g_i&=g_{\varepsilon_i,\pi_i}:\Delta_{\varepsilon_i}^{\ell_i}\longrightarrow\Delta_{\varepsilon_i}^{\operatorname{long}(\pi_i)}\ ; 
\end{aligned}
\]
$g_i$ lui-même étant représenté comme un $*$-produit de $\ell_i$ facteurs \inserted{indécomposables}, de sorte que \inserted{$f$ et} $g$ apparaissent sous forme
\scanpage{316}

\mspage{55}
% archive sheet number: 315

de $*$-produits de $n$, resp. de $\sum_{i\in[1,n]}\ell_i$ facteurs de la forme $f_{\varepsilon,\unsure{\nu|n}}$. Pour terminer de tenir en mains l'objet \unsure{$1/2$ simpl.} tronqué $(\operatorname{Nerf}_i(\underline{\mathrm{Ty}})\unsure{]|)}_{0\leq i\leq2}$, il reste à décrire les \inserted{deux} applications de dégénérescence $\operatorname{Nerf}_1\unsure{\text{ --- }|\longrightarrow}\operatorname{Nerf}_2$ (cas où $f$ ou $g$ est une flèche identique), et surtout l'application de \uline{composition} des flèches
\[
\operatorname{Nerf}_2(\underline{\mathrm{Ty}})\xrightarrow{\gamma}\operatorname{Nerf}_1(\underline{\mathrm{Ty}})=\operatorname{Fl}(\mathrm{Ty}).
\]
Il suffit de donner $\gamma$ sur les générateurs, et on a \unsure{tautologiquement}, par définition des $f_{\varepsilon,\ell,\pi}$
\[
\gamma(f_{\varepsilon,\ell,\pi})=f_{\varepsilon,\operatorname{long}(\pi)}.
\]
(NB Il n'existe \inserted{au plus} qu'un seul morphisme de $\Delta_\varepsilon$ i.e. (de $\Delta_+$ ou de $\Delta_-$) dans un intervalle $I_\tau$. Il existe sss $\tau$ est de la forme $\Delta_\varepsilon^\ell$, $\ell\in\mathbb N$.)

Le cas $g=\mathrm{id}$ correspond au cas où
\[
\pi_i(j)=1\qquad\forall j\in[1,\ell_i],
\]
les $f_{\varepsilon_i,\ell_i,\pi_i}$ sont \unsure{donc|alors} déterminés en fonction des $f_{\varepsilon_i,\ell_i}$, et l'application de dégénérescence est
\[
 f_{\deleted{\unsure{+}}\varepsilon,\ell}\longmapsto f_{\varepsilon,\ell,\underline{1}_{\deleted{\unsure{n}}\ell}}
\]
\margin{(où $1_\ell$ est la fonction \unsure{conste|const.} de val. $1$ sur $[1,\ell]$.)}
\deleted{\unsure{du}} \unsure{Cas|Le} / $f=\mathrm{id}_{\unsure{x|y}}$ \inserted{est celui où} les $\ell_i$ sont égaux à $1$. Alors \deleted{les \unsure{fonctions}} \deleted{$\unsure{\pi_i}$} chaque $\pi_i:[1,1]\to\mathbb N$ peut être vu comme un \deleted{\unsure{enti}} élément de $\mathbb N$, la donnée des $(\varepsilon_i,\ell_i,\pi_i)$ revient à celle des $(\varepsilon_i,n_i)\in[\pm1]\times\mathbb N$, i.e. la donnée \deleted{\unsure{du système}}
\scanpage{317}

\mspage{56}
% archive sheet number: 316

d'une suite de $n$ générateurs parmi les $f_{\varepsilon,\ell}$. L'application de dégénérescence $\operatorname{Nerf}_1\unsure{\text{ --- }|\longrightarrow}\operatorname{Nerf}_2$ est \unsure{ainsi|ici} donnée sur les générateurs par
\[
f_{\varepsilon,\ell}\longmapsto f_{\varepsilon,1,\ell}.
\]

Heureusement, on n'a plus à prouver l'associativité de la composition.

\uline{NB} La simplicité de ces descriptions \deleted{\unsure{me}} renforce mon sentiment que la catégorie $\underline{\mathrm{Ty}}$ définie à présent, avec les \uline{morphismes d'intervalles orientés croissants} $I_\tau\text{ --- }I_{\tau'}$ comme flèches de $\underline{\mathrm{Ty}}$, est bien ``la bonne''. On va déjà voir si elle donne une description sympa des catégories de fractions $\unsure{X.\Sigma^{-1}|X\Sigma^{-1}|Y.\Sigma^{-1}|Y\Sigma^{-1}}$, sans tortillements comme précédemment!

Mais on a promis de déterminer $\operatorname{Nerf}_r(\underline{\mathrm{Ty}})$, pour \unsure{tt|$\forall$} $r\in\mathbb N$. \deleted{L'élément \unsure{gé}} Le générateur général de $\operatorname{Nerf}_r$ est de la forme (\inserted{pour} $\varepsilon\in\{\pm1\}$, fixé)
\diagram{\[
\begin{tikzcd}[ampersand replacement=\&,column sep=small]
\underset{\substack{\Vert\\\Delta_\varepsilon^{\ell_0}\ \text{avec }\ell_0=1}}{\Delta_\varepsilon}\arrow[r,"f_1"] \&
\Delta_\varepsilon^{\ell_1}\arrow[r,"f_2"] \& \Delta_\varepsilon^{\ell_2}\arrow[r] \&
\cdots\arrow[r,dash,"f_r"] \& \Delta_\varepsilon^{\ell_r}
\end{tikzcd}
\tag*{\unsure{($\alpha$)|(3)|ds}}
\]}
où les $f_i$\inserted{$(i\in[1,r]):\Delta_\varepsilon^{\ell_{i-1}}\longrightarrow\Delta_\varepsilon^{\unsure{\ell_i|\ell_{i-1}}}$} sont donnés par \deleted{\unsure{des}} \deleted{\unsure{p}} applications
\[
\pi_{i-1}:[1,\ell_{i-1}]\to\mathbb N\qquad
\text{telle que}\quad
\underset{\substack{\Vert\,\mathrm{déf}\\
\sum_{j\in[1,\ell_i}\pi_i(j)}}{\operatorname{long}(\pi_{i-1})}
=\ell_{\unsure{i|i-1}}
\]
(N.B. $\pi_0:[1,1]=\{1\}\to\mathbb N$ est \unsure{identifié|identifiée} à $\operatorname{long}(\pi_0)=\pi_{\unsure{0|00}}(1)=\ell_1$)
\scanpage{318}

\mspage{57}
% archive sheet number: 317

Pour ne pas faire dépendre la définition du domaine des $\pi_i$ $(i\in[0,r-1])$ des entiers $\ell_i$, on peut aussi \unsure{considérer} les $\pi_i$ comme des applications
\[
\mathbb N^*\overset{\pi}{\unsure{\text{ --- }|\longrightarrow}}\mathbb N
\qquad(\mathbb N^*=\{i\in\mathbb N\mid i\geq1\}
\]
à supports finis (i.e. tels que $\pi(\unsure{n|i})=0$ pour $i$ grand), et exiger de façon précise
\[
\begin{aligned}
&\operatorname{supp}\pi_0\subset\{1\}\\
&\operatorname{supp}\pi_1\subset[1,\operatorname{long}(\pi_0)]\\
&\quad\vdots\\
&\boxed{\operatorname{supp}\pi_i\subset[1,\operatorname{long}(\pi_{i-1})]}
\quad\text{pour}\quad\boxed{i\in[0,r-1]}
\end{aligned}
\tag*{($**$)}
\]

D'autre part, \deleted{$\unsure{f_i}$} posant $\ell_i=\operatorname{long}(\pi_{\unsure{i-1|i}})$ (pour $i\in[1,\unsure{r|r-1}]$)
\[
 f_i:\Delta_\varepsilon^{\ell_{i-1}}\longrightarrow\Delta_\varepsilon^{\ell_i\,(=\operatorname{long}\pi_{i-1})}
\]
est défini comme
\[
 f_i=f_{\varepsilon,\pi_{i-1}(1)}*f_{\varepsilon,\pi_{i-1}(\unsure{2|1})}*
 \cdots*f_{\varepsilon,\pi_{i-1}(\ell_{i-1})},
\]
de sorte que l'élément $(\unsure{\alpha|*|x})$ de $\operatorname{Nerf}_r(\underline{\mathrm{Ty}})$ est déterminé de façon unique par la donnée des \deleted{$\unsure{\pi}$} $\varepsilon\in\{\pm1\}$ et de
\[
\pi_{\bullet}=(\pi_i)_{\deleted{\unsure{0\leq}}i\in[0,r-1]}
 \in\big(\operatorname{Hom}_{\mathrm{Ens}}(\mathbb N^*,\mathbb N)\big)^r
\]
satisfaisant les conditions $(**)$. \deleted{\unsure{D}} Soit $F_{\varepsilon,\pi_\bullet}$ \deleted{\unsure{le|l'élé}} cet élément de $\operatorname{Nerf}_r(\underline{\mathrm{Ty}})$, \deleted{Alors les}\inserted{et} \deleted{$\unsure{(F_{\varepsilon,\pi_\bullet})_{\varepsilon,\pi_\bullet}}$}\inserted{soit} $N_r$ l'ensemble des $\pi_\bullet$ qu'on vient d'expliciter, Alors
\[
(F_{\varepsilon,\pi_\bullet})_{(\varepsilon,\pi_\bullet)\in\{\pm1\}\times N_r}
\]
est un système de générateurs libre du monoïde $\operatorname{Nerf}_r(\underline{\mathrm{Ty}})$,

Si $\ell_i=0$\inserted{(i.e. $\boxed{\pi_{i-1}=0}$)}, on doit avoir \inserted{$\operatorname{supp}\pi_{\unsure{i+1|i-1}}=\varnothing$ i.e.} $\pi_{\unsure{i-1|i+1}}=0$, d'où $\ell_{i+1}=0$, et de $=$ pour les $\ell_j$ et $\pi_j$ suivants. Géométriquement, \deleted{\unsure{cette}}\inserted{l'hyp.} s'exprime par la condition que le but
\scanpage{319}

\mspage{58}
% archive sheet number: 318

de $f_i$ est l'élément $1$ de $\mathrm{Ty}$ (correspondant à l'intervalle ponctuel), et comme $\operatorname{Hom}(1,\tau)=\varnothing$ pour $\tau\neq1$, cela signifie que les flèches qui suivent $f_i$ doivent être toutes égales à la flèche $\mathrm{id}_1:1\unsure{\text{ --- }|\to}1$.

\uline{Remarque} Dans la description combinatoire de $\underline{\mathrm{Ty}}$ et de $\operatorname{Nerf}_{\unsure{\bullet|0}}(\underline{\mathrm{Ty}})$, l'ensemble des ``signatures'' $\{\pm1\}$ joue un rôle assez \unsure{bidon}. Il peut être remplacé par n'importe quel ensemble $\Sigma$ (dit ``ensemble des signatures''), et \deleted{\unsure{les}} on construit alors une catégorie monoïdale (ou un monoïde catégorique) ayant comme ensemble d'objets le monoïde libre engendré par $\Sigma$, interprétés comme des intervalles d'entiers $[0,n]$, munis de ``signatures'' $\in\Sigma$ pour chaque intervalle de transition, \inserted{ou ce qui revient au $=$, des ``graphes'' \deleted{\unsure{pointés}}\inserted{orientés} $\Sigma$-pondérés de type combinatoire un \deleted{\unsure{dist}} segment $\Delta_{\unsure{i|1}}$,} ce qui revient à la donnée d'une suite \inserted{(ordonnée)} de \deleted{\unsure{des}} éléments $\deleted{\unsure{\Sigma|E}}\varepsilon_1,\varepsilon_2,\unsure{\text{ --- }|\ldots},\varepsilon_n\in\Sigma$, i.e. d'un ``mot'' $\varepsilon_1\varepsilon_2\ldots\varepsilon_n$, qu'on préférera écrire $\Delta_{\varepsilon_1}\Delta_{\varepsilon_2}\ldots\Delta_{\varepsilon_n}$. On peut définir de façon évidente les hom. \uline{croissants} de tels graphes pondérés orientés (sous-entendu : avec l'élément source et l'élément but dans la structure, à respecter par les morphismes), sans avoir besoin d'une interprétation catégorique. La fonctorialité de l'opération de composition des intervalles pondérés, pour cette notion de morphisme, est évidente. \unsure{Le|La}
\scanpage{320}

\mspage{59}
% archive sheet number: 319

choix $\Sigma=\{\pm1\}$ est \deleted{\unsure{assurément|évidemment} \unsure{bien}} à ma connaissance le seul utile, à cause de son lien avec la notion de chemin et le calcul du $\pi_0$ d'un objet de $X$, etc.

7. \uline{Explicitation des catégories} \inserted{\uline{de fractions}} $X\Sigma^{-1}$ en termes \uline{des} $\underline{\mathrm{Ch}}(X;a,b)$.

Si $X$ est une catégorie, on désigne par $\underline{\mathrm{Ch}}(X)$ la catégorie des chemins de $\unsure{x|X}$ (cf \S{} $\deleted{\unsure{5}}\unsure{4|5}$), et pour $a,b\in\operatorname{Ob}X$, par
\[
\underline{\mathrm{Ch}}(X;a,b)\quad\text{ou simplement}\quad\underline{\mathrm{Ch}}(a,b)
\]
celle des chemins de $a$ vers $b$, i.e de source $a$, but $b$. L'\unsure{intention|intuition} géométrique (qu'on explicitera), c'est qu'un tel chemin \inserted{$c$} définit un chemin ``direct'' de $a$ vers $b$ (i.e. un chemin de type $\tau'$ \deleted{\unsure{p|pur}} ``direct'', i.e. \deleted{\unsure{qui}} toutes les flèches de transition sont directes, i.e. $\tau'$ de la forme $\Delta_+^n$), de $=$ longueur que le chemin initial $c$, obtenu en inversant toutes les flèches rétrogrades \unsure{dans} le chemin $c$. Cette opération est effectivement possible dans $X$ lui-même sss la totalité \unsure{flèches|flèche} sont \unsure{inversible|inversibles}, On obtient alors une suite \unsure{composables|composable} de flèches, dont le composé total \deleted{\unsure{est}} (qu'on peut appeler ``l'effet utile'' du \deleted{\unsure{chemin}})
\scanpage{321}

\mspage{60}
% archive sheet number: 320

est une flèche de $a$ vers $b$. Dans le cas général, les chemins de $a\unsure{)|,}$ vers $b\unsure{)}$ peuvent être vus comme une sorte de ``flèches généralisées'' allant de $a$ vers $b$, dont la composition \deleted{est la} serait la composition \unsure{évidente}, et les identités les ``chemins vides'' relatifs aux $a\in\operatorname{Ob}X$. \deleted{\unsure{D|Il}}\unsure{On construit|On obtient} d'ailleurs bien ainsi une catégorie, ayant $=$s objets que $X$, mais où les
\[
\operatorname{Hom}_{\mathrm{ch}}(a,b)=\mathrm{Ch}(a,b)=\operatorname{Ob}\underline{\mathrm{Ch}}(a,b)
\]
sont prohibitivement gros. \deleted{\unsure{Il}} Cela s'exprime \inserted{notamment} par le fait que les seules flèches inversibles dans $X_{\mathrm{ch}}$ sont les identités! On n'a pas de foncteur naturel entre $X_{\mathrm{ch}}$ et $X$, \deleted{\unsure{dans le}}\inserted{qui} soit l'identité sur les objets. Car à un chemin $a\overset{c}{\unsure{\text{ --- }|\longrightarrow}}b$ on ne peut associer de $X$-morphisme $a\text{ --- }b$ (sauf si $c$ de type $\Delta_+$), et à un $X$-morphisme $a\overset{u}{\text{ --- }}b$ on peut associer un chemin \inserted{$c(u)$} de type $\Delta_+$, mais \deleted{\unsure{ce}} \inserted{si $a\overset{u}{\text{ --- }}b\overset{v}{\text{ --- }}c$,} $c(vu)$ est \inserted{encore} de type $\Delta_+$ alors que $c(v)c(u)$ est de type $\Delta_+^2$, ils ne sont pas égaux. Par contre, il y a une flèche \inserted{naturelle} entre ces deux \deleted{\unsure{chemins}} éléments de $\underline{\mathrm{Ch}}(a,\unsure{c|b})$, flèche au dessus de \deleted{l'\unsure{hom. entre}} l'unique hom.
\[
\Delta_+\xrightarrow{f_{+,2}}\Delta_+^2
\]
\scanpage{322}

\mspage{61}
% archive sheet number: 321

On a alors
\[
f_{+2}^{*}\bigl(\underbrace{a\overset{u}{\text{ --- }}b\overset{v}{\text{ --- }}c}_{c(v)c(u)}\bigr)
=\underset{c(vu)}{(a\xrightarrow{vu}c)}
\]
d'où une \unsure{flèche|flèches} canonique
\[
c(vu)\longrightarrow c(v)c(u)\qquad\text{dans }\underline{\mathrm{Ch}}(a,c)
\]

Cela montre que $c(vu)$ et $c(v)c(u)$ sont \unsure{donc} au moins dans une même composante connexe de $\underline{\mathrm{Ch}}(a,c)$, ou \deleted{\unsure{encore}} (comme on dit aussi) qu'ils sont ``homotopes''. Cela suggère \inserted{\unsure{ici}} de \deleted{\unsure{déf}} modifier la définition des $\operatorname{Hom}_{\mathrm{ch}}(a,b)$ en
\[
\pi(a,b)=\pi_0\bigl(\underline{\mathrm{Ch}}(a,b)\bigr).
\]

Le foncteur de composition des chemins
% slanted
\margin{Ce foncteur est déduit des foncteurs $\underline{\mathrm{Ch}}_\tau(a,b)\times\underline{\mathrm{Ch}}_{\tau'}(b,c)\longrightarrow\underline{\mathrm{Ch}}_{\tau\tau'}(a,c)$ \unsure{evidents|évidents}, cf \unsure{pag|pge} \unsure{7.7|77|47}}
\[
\begin{aligned}
\underline{\mathrm{Ch}}(a,b)\times\underline{\mathrm{Ch}}(b,c)&\;\text{ --- }\;\underline{\mathrm{Ch}}(a,c)\\
(\unsure{a|c},c')&\;\unsure{\mathrel{\mapstochar\relbar}|\longmapsto}\;c'\circ c
\end{aligned}
\]
induit une \deleted{\unsure{``compos}} application sur les $\pi_0$, ce qui donne une loi de composition
\[
\pi(a,b)\times\pi(b,c)\;\text{ --- }\;\pi(a,c)
\]
sur les flèches. Il est immédiat qu'elle est associative (la composition des chemins l'étant), et unitaire. Elle admet comme flèche identique en $x$ soit le \deleted{\unsure{mot}} \deleted{chemin} ``le mot vide en $x$'', lequel est \inserted{(évidemment)} homotope dans $\underline{\mathrm{Ch}}(x,x)$ au \unsure{chemin} ``dégénéré'' $x\xrightarrow{\mathrm{id}}x$
